New sequencing technologies have fundamentally expanded our ability to study processes across the central dogma, from the genetic architecture of individuals \citep{green2010draft, gravel2011demographic, tennessen2012evolution, lazaridis2014ancient, haak2015massive} to changes in gene expression over time at the single-cell level \citep{shapiro2013single, grun2015design} during cells' life cycle. These diverse measurements offer complementary perspectives on biological processes such as human history, cancer development, and immune system function. At the same time, the cost of sequencing has decreased substantially \citep{christensen2015assessing, park2016trends, schwarze2018whole, nawy2014single, ding2022temporal, potter2018single}, enabling increasingly large-scale studies. With these data, we can answer more scientific questions than before. Many scientific questions take the form of forecasting---for example, predicting the number of unique genetic variants in large, costly follow-up studies based on a small pilot sample, or forecasting changes in gene expression over time during immune cell life cycles.

Forecasting is fundamentally challenging, as it requires making predictions beyond the support of observed data, such as extrapolating to new sample sizes or future time points. The inherent noisiness of genomics data further complicates this task. Nevertheless, biologists often possess imperfect but valuable domain knowledge that can inform forecasting. An ideal method should leverage such knowledge while remaining flexible to both the data and the uncertainty or misspecification in that knowledge. More broadly, these challenges reflect a need of combining information from different sources, necessitating methods that are both flexible enough to fit the data and structured enough to incorporate potentially imperfect domain knowledge. 


In this thesis, we address two forecasting tasks that instantiate these challenges. The first arises naturally from large-scale population genetics studies targeting novel (rare) genetic variants \citep{lek2016analysis, karczewski2020mutational}. These variants play a central role in many diseases and represent potential targets for new therapeutics \citep{szustakowski2021advancing}. In this setting, a forecasting method must estimate the number of novel variants in a new (follow-up) study across multiple populations (e.g., different cancer types), given data from a typically small pilot study—that is, forecasting beyond the observed \textit{sample sizes}. It is also of interest to predict related summaries, such as the number of shared or population-specific (private) variants. In all cases, it is crucial to account for variants shared across groups as well as those unique to individual populations. Limited pilot data may be insufficient to reliably estimate quantities such as the upper bound on the number of variants, while multi-population settings require estimating correlation structures from sparse observations.

To address these challenges, we develop Bayesian nonparametric models for predicting the number and structure of unseen genetic variants across multiple populations, accounting for both shared and population-specific variation. The Bayesian nonparametric approach avoids the need to specify a fixed upper bound on the number of variants, thereby maintaining flexibility, while still allowing the incorporation of domain knowledge such as power-law growth patterns \citep{Masoero2022}. We further adopt an empirical Bayes-like strategy to allow the data to inform these structural assumptions.

The second task arises from single-cell sequencing studies of gene expression dynamics over a cell's life cycle, measured through mRNA concentration. One can model these dynamics using e.g., stochastic differential equations \citep[][SDEs]{pratapa2020benchmarking}. Although forecasting has been extensively studied when a single trajectory is observed \textit{over time}, these techniques cannot be used in the single cell sequencing setup. Because in practice measuring mRNA concentration requires destroying the cell, it is impossible to track individual cells' gene expression level longitudinally to use classic timeseries techniques. The goal is therefore twofold: (1) to reconstruct the distribution of unobserved trajectories and (2) to forecast trajectories forward in time. In this setting, learning an SDE from snapshot data is fundamentally ill-posed: without additional assumptions, infinitely many SDEs can produce the same observed snapshots \citep{lavenant2024toward}.

To address these challenges, we develop methods for inferring latent stochastic dynamics from snapshot observations and forecasting their evolution over time. Our approach accommodates a large and flexible class of SDEs (e.g., using nonparametric function classes or neural networks) while incorporating domain knowledge in a form that does not require it to be encoded as a single fully specified model. Similar to an empirical Bayes framework, our method allows prior knowledge to be adapted based on observed data.

Although these problems arise in different domains, they are unified by a common goal: enabling reliable forecasting under limited and indirect observations. Addressing this goal requires probabilistic methods that are flexible enough to adapt to data while structured enough to incorporate domain knowledge, without requiring that such knowledge be complete or perfectly specified.
